\documentclass[11pt,a4paper]{article}
\usepackage[margin=22mm]{geometry}
\usepackage{amsmath,amssymb}
\usepackage{booktabs}
\usepackage{graphicx}
\usepackage{microtype}
\usepackage{float}
\usepackage[hidelinks]{hyperref}
\setlength{\parindent}{0pt}
\setlength{\parskip}{0.55em}
\graphicspath{{../play/out/research/}}

\title{Three Free Slots or Two?\\\large Why they looked the same, and where they differ}
\author{Minh Quan Nguyen}
\date{October 2026}

\begin{document}
\maketitle

\begin{abstract}
In our earlier experiments a scheme with three free slots (one overlapping both others) performed exactly
like a scheme with two. We show why: with every customer drawing full power for the whole slot, an
overlapping slot can never lower a circuit's peak, and our model asked a different question from the
proposal (doc/references/New.md), in which almost everyone keeps the mandated 11:00--14:00 slot and only
``heavy hitters'' move. Rebuilding the model around the proposal's premise (a phase can host one 40 kW
heavy hitter on a cloudy day, but not two) and adding the objectives the proposal cares about, the
three-slot scheme does differ. It is a dial between disruption and network safety: moving only heavy
hitters (0.9\% of customers) raises the share of phases needing no upgrade from 83\% to 92\%, and moving
10\% of customers raises it to 96\%, against 99.5\% for moving everyone into two slots (32\%). It keeps
heavy users' full benefit, unlike a power cap, and it spreads switching over more times, cutting the
largest step in Victorian load by 17\% against two slots and 57\% against the policy at 50\% uptake.
\end{abstract}

\section{Why two and three slots looked the same}

\begin{enumerate}\itemsep2pt
\item \textbf{Structure.} With the middle slot overlapping both others, the day splits into segments
  carrying \(m_1\), \(m_1+m_2\), \(m_2+m_3\) and \(m_3\) customers' extra load. The peak
  \(\max(m_1+m_2,\,m_2+m_3)\) is never below \(\max(m_1,m_3)\), which two slots already achieve, so a
  customer placed in the middle can only raise it. In all 861 circuit compositions of the report's model
  the lowest peak with three slots equals that with two (Figure~\ref{fig:r1}).
\item \textbf{Framing.} Every medium and large customer could be moved, so ``two slots'' meant moving
  31\% of customers away from the mandated window. Disruption was never measured, yet it is the
  proposal's main argument.
\item \textbf{Calibration.} In the report's model the medium users alone overload a shared slot
  (800 A of base load plus about 240 A for twelve medium users, against a 1000 A limit). The proposal's
  premise, that a circuit can host one heavy hitter but not two, did not hold, so moving heavy hitters
  could not help and only moving everyone did.
\item \textbf{Soak.} Soak depends on which hours the slots cover. The middle slot adds no hours, and
  spreading load never loses soak, so three slots could at best tie.
\item \textbf{Load timing.} With All-On, or with random start times, use in the middle slot overlaps use
  in both others, so staggered slot openings cannot help.
\item \textbf{Missing objectives.} Disruption, the benefit kept by heavy users, and the shape of the
  system load (peak and switching steps), where the middle slot plausibly matters, were not measured.
\end{enumerate}

\begin{figure}[H]\centering
\includegraphics[width=0.45\textwidth]{r1_diagnosis.png}
\caption{Lowest circuit peak with two slots (10--13, 13--16) and three slots (adding 11--14), for every
composition of the report's model. All points lie on the diagonal.}
\label{fig:r1}
\end{figure}

\section{What the three-slot scheme is for}
The proposal keeps the mandated slot 11:00--14:00 for almost everyone and adds two non-overlapping relief
slots, 9:00--12:00 and 12:00--15:00, for heavy hitters (for example two level-2 EV chargers, 40 kW). When
two heavy hitters share a phase they move to the two relief slots, so the phase never carries both at
once. The middle slot is not there to add network headroom; it keeps the tariff simple and keeps most
load at peak solar. The comparison that matters is therefore the policy against the proposal, and
against the alternatives a regulator would consider: moving everyone into two slots, or capping free
power.

\section{Model}
\paragraph{Phase.} A circuit is one phase of a distribution line with \(N\) customers. On a cloudy day
every customer draws a base current \(I_0\); moderate free-window users add \(\Delta I_M\) inside their
slot and heavy hitters add \(\Delta I_H\). Customer types are random (trinomial, shares \(p_M\) and
\(x\)). The phase limit is set so that an average phase can host one heavy hitter but not two:
\(T=N I_0+p_M N\Delta I_M+1.25\,\Delta I_H\). The values are illustrative and would be replaced by DNSP
data.

\begin{center}\small
\begin{tabular}{lll}
\toprule
Parameter & Value & Note\\
\midrule
\(N\) & 30 & customers per phase\\
\(I_0\) & 5 A & 1.2 kW base load on a cloudy midday\\
\(\Delta I_M\), \(p_M\) & 15 A, 30\% & moderate users (3.6 kW, e.g.\ hot water)\\
\(\Delta I_H\), \(x\) & 167 A, 2\% & heavy hitters (40 kW)\\
\(T\) & 493 A & base 150 + moderate 135 + 1.25 heavy hitters\\
Uptake \(u\) & 50\% & share of Victorian phases on the plan\\
Solar & 2025 VIC rooftop & seasonal weights as in the main report\\
\bottomrule
\end{tabular}
\end{center}

\paragraph{Strategies.}
\begin{itemize}\itemsep0pt
\item \textbf{Policy:} everyone in 11--14.
\item \textbf{Proposal:} everyone in 11--14; a phase with two or more heavy hitters sends them to 9--12
  and 12--15.
\item \textbf{Three slots, middle share \(a\):} the proposal, plus a share \(1-a\) of moderate users
  moved to the relief slots (\(a=1\) is the proposal; \(a=1/3\) splits users evenly).
\item \textbf{Two-slot split:} every free-window user in 9--12 or 12--15, lowest peak.
\item \textbf{Cap:} everyone in 11--14, free power capped at 7 kW per customer.
\end{itemize}

\paragraph{Metrics.} The probability that a phase needs no upgrade (its peak stays under \(T\)); the
share of customers moved out of 11--14; rooftop solar soaked; and, for Victoria (\(u\,N_{\mathrm{VIC}}/N\)
phases), the rise in the mean daily peak of operational demand plus the added load, the largest
5-minute step in the added load, and the free energy each heavy hitter receives.

\section{Results}

\subsection{Safety against disruption}

\begin{center}\footnotesize
\begin{tabular}{lcccccc}
\toprule
Strategy & No upgrade & Moved & Soak & Peak rise & Largest step & Free energy\\
& & & (MWh/day) & (MW) & (MW) & per heavy hitter\\
\midrule
Policy                 & 83.1\% & 0\%    & 6{,}523 & 598 & 2{,}726 & 120 kWh\\
Proposal               & 91.8\% & 0.9\%  & 7{,}054 & 428 & 2{,}212 & 120 kWh\\
Three slots, 1/3 each  & 97.8\% & 20.9\% & 7{,}690 & 201 & 1{,}168 & 120 kWh\\
Two-slot split         & 99.5\% & 32.0\% & 7{,}704 & 212 & 1{,}410 & 120 kWh\\
Cap, 7 kW              & 100\%  & 0\%    & 5{,}001 & 132 & 1{,}769 & 21 kWh\\
\bottomrule
\end{tabular}
\end{center}
Heavy hitters 2\%, \(N=30\), uptake 50\%.

\begin{figure}[H]\centering
\includegraphics[width=0.85\textwidth]{r7_frontier.png}
\caption{Share of phases needing no upgrade against the share of customers moved out of 11--14. The line
traces the three-slot scheme from the proposal (middle share 100\%) to all moderate users in the relief
slots (0\%). It uses a simple rule (moderate users split evenly, a single heavy hitter stays in 11--14),
so its end point is below the two-slot split, which chooses each phase's lowest-peak split.}
\end{figure}

\begin{enumerate}\itemsep0pt
\item \textbf{The proposal recovers half of the available gain while moving under 1\% of customers}
  (83\% to 92\% of phases safe, against 99.5\% for the two-slot split, which moves 32\%).
\item \textbf{The middle share is a dial.} Moving 7\% of customers gives 95\%, 10\% gives 96\%, 16\%
  gives 97\%; beyond that the gain flattens. Two slots offer only the end point.
\item \textbf{A cap is safest but defeats the benefit-sharing goal.} It keeps every phase safe without
  moving anyone, but heavy users get 21 instead of 120 kWh of free energy a day and Victoria soaks 23\%
  less solar than under the policy.
\end{enumerate}

\subsection{Probability of avoiding an upgrade}
\begin{figure}[H]\centering
\includegraphics[width=\textwidth]{r3_heatmaps.png}
\caption{Probability that a phase needs no upgrade, against the share of heavy hitters and the customers
per phase. Right: proposal with smart-charging timers, sessions of 1 hour from the slot opening and a
third heavy hitter placed in 11--14.}
\end{figure}
For 30 customers per phase and 2\% heavy hitters, 83\% of phases need no upgrade under the policy and
92\% under the proposal; at 6\% heavy hitters, 40\% and 64\%. The proposal's advantage grows with the
number of heavy hitters, as the reference note's \(p_2/p_3\) argument predicts.

\subsection{Smart-charging timers}
If heavy hitters charge for \(d\) hours from the opening of their slot, a session in 9--12 ends before
the moderate users' 11:00 start, and with \(d\le1\) h a third heavy hitter fits in 11--12 between the
morning and afternoon sessions.
\begin{figure}[H]\centering
\includegraphics[width=0.75\textwidth]{r4_timers.png}
\caption{Phase of 30 customers: probability of needing no upgrade against the share of heavy hitters, for
full-slot use (All-On) and for sessions of 1 and 1.5 hours from the slot opening.}
\end{figure}
\begin{enumerate}\itemsep0pt
\item \textbf{Short sessions help the proposal most}: at 6\% heavy hitters, 78\% of phases are safe with
  1-hour sessions against 64\% with full-slot use, because morning sessions avoid the moderate users.
\item \textbf{A third heavy hitter in the middle slot adds a further 7 points} (85\%) when sessions last
  at most 1 hour. With 1.5-hour sessions it overlaps the afternoon session and is worse than leaving it
  out (64\% against 78\%). The middle slot raises capacity only if sessions are shorter than the 1-hour gap between
  the 11:00 and 12:00 openings.
\end{enumerate}

\subsection{System level}
\begin{figure}[H]\centering
\includegraphics[width=\textwidth]{r5_system.png}
\caption{Victoria, against uptake: rise in the mean daily peak of operational demand plus the added load,
largest 5-minute step in the added load, and rooftop solar soaked. Adaptive: the middle share with the
lowest peak at each uptake (labels).}
\end{figure}
\begin{figure}[H]\centering
\includegraphics[width=0.7\textwidth]{r6_shapes.png}
\caption{Added Victorian load at 50\% uptake.}
\end{figure}
\begin{enumerate}\itemsep0pt
\item \textbf{Switching steps.} Under the policy every plan customer switches on at 11:00 and off at
  14:00, a 2.7 GW step at 50\% uptake. Two slots halve it (1.4 GW); three slots with users split evenly
  cut it to 1.2 GW (\(-17\%\) against two slots), because load switches at five times instead of three.
\item \textbf{Peak.} Up to 25\% uptake the added load stays below the evening peak and every strategy
  raises the daily peak by under 65 MW. Above that, concentrating everyone in 11--14 creates a new midday
  peak (+598 MW at 50\%, +3.2 GW at 100\%); spreading users over the relief slots keeps it to 35--40\% of that;
  among the strategies that keep heavy users' full benefit the adaptive three-slot share is lowest
  (+156 MW at 50\%), and only the cap is lower (+132 MW).
\item \textbf{Soak.} Spreading load also soaks more rooftop solar at high uptake (7{,}704 against
  6{,}523 MWh/day at 50\%); the proposal alone gains 8\%, and the cap loses 23\%.
\end{enumerate}

\section{Scorecard}
\begin{center}\footnotesize
\begin{tabular}{lccccc}
\toprule
& Network & Disruption & Heavy users' & Solar & Switching\\
& safety & & benefit & soak & steps\\
\midrule
Policy           & low      & none   & full    & baseline & largest\\
Proposal         & medium   & tiny   & full    & +8\%   & large\\
Three slots (dial) & medium--high & chosen & full & +18\% & smallest\\
Two-slot split   & highest  & large  & full    & +18\%  & medium\\
Cap              & highest  & none   & cut 82\% & \(-23\%\) & medium\\
\bottomrule
\end{tabular}
\end{center}
Changes in soak against the policy at 50\% uptake.

\section{What this means}
\begin{itemize}\itemsep0pt
\item Three slots do not beat two on circuit headroom; the overlap makes that impossible. They win where
  the proposal's goals are: few customers disturbed, heavy users keep their benefit, and the operator can
  trade disruption for safety and smooth the system load by choosing the middle share.
\item The strongest case to a DNSP is the frontier: 92\% of phases safe for 0.9\% of customers moved,
  96\% for 10\%, without rewiring or upgrades.
\item The strongest case to the market operator is the switching step: a single mandated slot creates a
  multi-gigawatt step at 11:00 and 14:00 at high uptake; three slots spread it.
\end{itemize}

\section{Limitations and data needed}
The phase parameters are illustrative; DNSP data are needed on customers per phase, phase ratings,
cloudy-day midday demand, and the number and peak power of heavy hitters. Moderate users draw for their
whole slot; heavy-hitter timers are assumed to start at the slot opening; customers are independent.
Switching steps assume synchronised timers; diversity in real starts would smooth them. Seasonal weights
scale only the load counted for soak.

\end{document}
